\section{Discussion}
\label{sec:discussion}

\paragraph{Why do structured tasks saturate?}
We hypothesise that structured-output quality is governed primarily by
\emph{schema compliance}, the model's ability to produce a response in the
required format (e.g., a sentiment label or a four-field JSON object).
Schema compliance is a near-discrete capability: the model either produces
the correct format or it does not.
Once prompt elaboration crosses the threshold that activates this capability,
additional instruction tokens provide no further signal.
This is consistent with the sigmoid shape of classification and product
extraction curves: a steep rise as the model ``activates'' the schema,
followed by a plateau.

\paragraph{Why do open-ended tasks not saturate?}
Open-ended tasks require the model to retrieve factual knowledge (QA) or
execute multi-step reasoning (math).
These capabilities are bounded by what was learned during pre-training, not by
how clearly the task is described.
Adding more instruction tokens cannot supply knowledge that was never acquired,
nor can it increase the model's effective reasoning depth.
At level 1 (bare input), capable models already provide near-correct answers;
weaker models fail not because they misunderstood the task, but because they
lack the requisite knowledge or reasoning.
In both cases, further elaboration is uninformative.

\paragraph{Practical guidelines.}
Our saturation points use a 95\%-of-asymptote threshold, making them
conservative estimates.
Practitioners can use these as upper bounds for prompt budget allocation:

\begin{itemize}
  \item \textbf{Sentiment classification}: a concise task description with
    label definitions ($\approx$50 tokens) is sufficient for most frontier
    models; adding a worked example ($\approx$300+ tokens) provides no benefit.
  \item \textbf{Structured extraction}: models vary substantially (92--536
    tokens). Weaker or smaller models may require more schema detail;
    stronger models can infer the schema from minimal cues.
  \item \textbf{QA and math reasoning}: prompt length does not predict
    quality. Invest instead in retrieval-augmented generation (for QA) or
    chain-of-thought formatting (for reasoning), which improve quality through
    mechanisms other than instruction elaboration.
\end{itemize}

\paragraph{Model capability modulates saturation.}
Within classification, stronger models (gemini-flash, gpt-4o-mini) saturate
at lower token counts ($\approx$42--50 tokens) than the 70B Llama model (64
tokens).
The 8B Llama model fails to produce a significant curve at all, suggesting it
may require even more elaboration, or that its classification quality is too
noisy for a clean saturation signal.
These gradations suggest that saturation points could serve as a
\emph{diagnostic} for model instruction-following sensitivity.

\subsection{Limitations}
\label{sec:limitations}

\paragraph{Statistical power.}
With only $n = 7$ data points per curve, the F-test has limited power.
Some pairs with genuine saturation may be missed; the low-significance
borderline cases (classification for llama-3.1-8b at $p = 0.079$;
product extraction for kimi-k2 at $p = 0.093$) are likely true positives
that would reach significance with more levels.
Future work should use finer-grained level spacing (e.g., 15--20 levels).

\paragraph{Single judge model.}
All quality scores are produced by \texttt{gpt-4o-mini}.
Although we validate against heuristic metrics ($r = 0.986$), we do not
validate against human raters.
Judge bias favouring certain response styles could systematically inflate or
deflate quality estimates, though the cross-evaluator agreement suggests this
is minimal.

\paragraph{English only.}
All prompts and examples are in English.
Saturation behaviour may differ for low-resource languages where prompt
elaboration is more necessary to disambiguate tasks.

\paragraph{Level ordering.}
Our results reflect a specific ordering of elaboration layers (task label
before format before definitions before persona before example).
A different ordering might produce different saturation points.
Studying the effect of layer permutation is left for future work.
